\documentclass[border=10pt]{standalone}

% пакеты для математики
\usepackage{amsmath,amsfonts,amssymb,amsthm,mathtools}

\usepackage[british,russian]{babel} % выбор языка для документа

% Математические шрифты
\usepackage{unicode-math}
\setmathfont[math-style=upright]{[Neo Euler.otf]}

\usepackage{tikz}
\usepackage{xcolor}

\definecolor{blue}{RGB}{0,114,178}

\usetikzlibrary{calc,shapes,positioning}

\begin{document}
    \begin{tikzpicture}[line cap=round,line join=round,x=1.6cm,y=1.6cm]
        % ступеньки: высота равна значению функции в середине отрезка
        \foreach \i in {0,...,11}{
            \pgfmathsetmacro\a{\i*0.5}
            \pgfmathsetmacro\m{\a+0.25}
            \pgfmathsetmacro\b{1.5+0.7*sin(1.1*\m r)+0.35*sin(2.9*\m r)}
            \draw [line width=0.8pt,fill=blue!20,draw=blue!70!black] (\a,0) rectangle (\a+0.5,\b);
        }

        % оси
        \draw [->, line width=1.pt] (-0.3,0) -- (6.5,0) node[below] {\small $x$};
        \draw [->, line width=1.pt] (0,-0.3) -- (0,2.9) node[left] {\small $f(x)$};

        % сама функция
        \draw [line width=1.5pt,smooth,samples=100,domain=0:6.2] plot (\x,{1.5+0.7*sin(1.1*\x r)+0.35*sin(2.9*\x r)});

        % одна ступенька крупным планом
        \draw [line width=1.pt] (3,-0.07) -- (3,0.07);
        \draw [line width=1.pt] (3.5,-0.07) -- (3.5,0.07);
        \draw (2.93,0) node[below] {\small $a_i$};
        \draw (3.75,0) node[below] {\small $a_i + h$};
    \end{tikzpicture}
\end{document}
